\documentclass[10pt,a4paper]{amsart}
\usepackage[margin=19mm]{geometry}
\usepackage{amsmath,amssymb,mathtools}
\usepackage[scr=rsfs,cal=euler]{mathalfa}
\usepackage{xfrac}
\usepackage[unicode,hidelinks,bookmarks=false]{hyperref}
\newcommand{\ie}{i.e.\ }
\newcommand{\cf}{cf.\ }
\newcommand{\resp}{respectively\ }
\newcommand{\Subsection}{\subsection}
\providecommand{\nobreakdash}{\nobreak\hbox{-}}
\setlength{\emergencystretch}{2em}
\multlinegap0pt
\usepackage{tikz}
\usepackage{amssymb}
\usepackage{mathtools}
\newtheorem{lem}{Lemma}
\DeclareMathOperator{\GL}{GL}
\def\R{\mathbb{R}}
\def\x{\times}
\title{Fourier transforms and a filtration on the Lagrangian cobordism group of tori}
\author{\'Alvaro Mu\~niz-Brea}
\date{Source: arXiv:2411.16543v2 (CC BY 4.0)}
\begin{document}
\maketitle

\noindent\emph{Source and licence.} Fourier transforms and a filtration on the Lagrangian cobordism group of tori. Version arXiv:2411.16543v2 (CC BY 4.0), distributed by arXiv under the Creative Commons Attribution 4.0 International licence, http://creativecommons.org/licenses/by/4.0/, as recorded in the arXiv licence metadata for this exact version and shown on the abstract page https://arxiv.org/abs/2411.16543v2. Full licence notice: \texttt{licenses/arxiv-2411.16543-CC-BY-4.0.txt}.\par\smallskip

\noindent\textit{Editorial scope.} Counted unit: the article's lem lem:perturbations\_surject (source0.tex lines 1315--1318). The article's complete original proof of it is delivered with it (source0.tex lines 1319--1341, one \texttt{proof} environment of 23 lines). No mathematics is written, repaired or re-derived: the statement and the proof are the article's own lines, in the article's own notation, and no retained line is substituted or re-flowed.  No further source-local context is needed: every cross-reference of every retained line resolves inside the two delivered spans.  Nothing was omitted from any retained span, so the package carries no omission record.  No retained line contains a citation, so the wrapper carries no bibliography and no bibliography entry.  No retained line contains a figure environment, an image-inclusion command or drawing code, so no image is delivered and the package declares no asset dependency.  Editorial formatting only, in addition: an AMS \texttt{amsart} preamble that restates the article's own declarations and macros used by the retained lines, namely the environments \texttt{lem} on separate counters, together with the standard packages those lines need.  The retained environments print in this document as Lemma 1 in order of appearance, whereas the article prints the same items with the numbering of its own full document.  The complete native source of the article as distributed in the arXiv e-print for this exact version is bundled unchanged as \texttt{sources/169074/source0.tex}.
\begin{lem}\label{lem:perturbations_surject}
    Let $V(f)$ be a regular tropical hypersurface and $S \subset V(f)$ any cell.
    Then the map $\Delta_S$ is surjective onto a neighborhood of the origin.
\end{lem}

\begin{proof}
    Let $p \in V(f)^{(n-1-k)}$ be a point in the $(n-1-k)$-stratum of $V(f)$.
    If $\psi: U \subset B \to \R^n$ is an affine chart around $p$, then by regularity of $V(f)$ there exists $A \in \GL(n;\R)$ such that
    $$
    (A \circ \psi)(V(f)) = \mathcal P_{k} \x \R^{n-1-k},\quad (A\circ\psi)(p)=0
    $$
    where $\mathcal P_{k} \subset \R^{k+1}$ is the $k$-dimensional tropical pair of pants.
    (Note that a priori $A$ has real coefficients, hence $A \circ \psi$ is not an affine chart, but we do not need that for the proof.)
    It follows that the map $\Delta_S$ is surjective onto a neighborhood of the origin if and only if the map 
    $$
        \Delta_{\{0\}}: \R^{k+2} \to T_0 \R^{k+1} \cong \R^{k+1}
    $$
    associated to $\mathcal P_{k}$ is surjective. 
    We have that $\mathcal P_{k} = V(h_0)$ for
    $$
        h_\delta = \max\{\delta_0,x_1+\delta_1,\dots,x_{k+1}+\delta_{k+1}\}.
    $$
    A simple computation shows that the map $\Delta_{\{0\}}$ is given by
    $$
        (\delta_0,\dots,\delta_{k+1}) \mapsto (\delta_1-\delta_0,\dots,\delta_{k+1}-\delta_0)
    $$
    which is clearly surjective.
\end{proof}

\end{document}
